\documentclass[11pt]{article}
\newcommand{\shorttitle}{Maximum principal nullity over the binary field}
\usepackage[letterpaper,left=1.04in,right=1.04in,top=.92in,bottom=.9in,headheight=14pt]{geometry}
\usepackage[T1]{fontenc}
\usepackage{lmodern,amsmath,amssymb,amsthm,mathtools,booktabs,microtype,xcolor,fancyhdr,hyperref}
\definecolor{ink}{HTML}{173345}
\definecolor{accent}{HTML}{235E76}
\hypersetup{colorlinks=true,urlcolor=accent,linkcolor=accent,citecolor=accent}
\setlength{\emergencystretch}{1.5em}
\pagestyle{fancy}\fancyhf{}
\fancyhead[L]{\small\sffamily\color{accent}\shorttitle}
\fancyhead[R]{\small\sffamily\color{accent}26 November 2026}
\fancyfoot[C]{\small\thepage}
\renewcommand{\headrulewidth}{.3pt}
\newcommand{\papertitle}[2]{\thispagestyle{plain}\noindent{\small\sffamily\color{accent}#1 \hfill 26 November 2026}\par\vspace{.65em}{\raggedright\color{ink}\fontsize{24}{28}\selectfont\bfseries #2\par}\vspace{.65em}}
\newtheorem{theorem}{Theorem}
\newtheorem{lemma}[theorem]{Lemma}
\newtheorem{corollary}[theorem]{Corollary}
\newtheorem{proposition}[theorem]{Proposition}
\theoremstyle{remark}\newtheorem*{remark}{Remark}
\newcommand{\F}{\mathbb F}
\newcommand{\E}{\mathbb E}
\newcommand{\Prob}{\mathbb P}
\newcommand{\Span}{\operatorname{span}}
\hypersetup{pdftitle={The largest principal kernel of a random binary orthogonal matrix}}
\begin{document}
\papertitle{Finite fields and probability}{The largest principal kernel of a random binary orthogonal matrix}
\noindent\textbf{Abstract.}
Choose $A$ uniformly from the matrices over $\mathbb F_2$ satisfying $A^{\mathsf T}A=I$. Among its $2^n$ principal submatrices, the largest kernel dimension is concentrated on the two integers $h-1,h$, where $h$ is the nearest integer to $\sqrt n$. At dimensions whose distance from the squares tends to infinity it equals $\lfloor\sqrt n\rfloor$ with probability tending to one. Both candidate probabilities have positive lower limits when $n=k^2-1$. We prove these statements by comparing the images of a few vectors with independent random vectors under row restrictions, and by controlling pairs of overlapping supports. This answers the typical-size question posed by Vaas and sharpens the upper bound in the linked ProveIt report. The result is unconditional.

\section{A threshold at \texorpdfstring{$r^2=n$}{r squared equals n}}
All ranks and nullities are over $\mathbb F_2$. For $S\subseteq[n]=\{1,\ldots,n\}$, $A[S,S]$ retains precisely the rows and columns indexed by $S$. Put
\[
 O(n,2)=\{A\in\mathrm{GL}_n(\mathbb F_2):A^{\mathsf T}A=I\},\qquad
 \nu(A)=\max_{S\subseteq[n]}\dim\ker A[S,S].
\]
The notation $O(n,2)$ here specifies the isometries of the dot product. The empty matrix has nullity zero.

\begin{theorem}[Threshold and two-point concentration]\label{thm:threshold}
There are absolute positive constants $c,C$ such that, for all sufficiently large $n$ and integers $\sqrt n/2\le r\le2\sqrt n$,
\begin{align}
 \Pr(\nu(A)<r)&\le C(2^{r^2-n}+2^{-r}),\label{eq:fail}\\
 c\min(1,2^{n-r^2})\ \le\ \Pr(\nu(A)\ge r)
       &\le\min(1,C2^{n-r^2}).\label{eq:hit}
\end{align}
If $h$ is the nearest integer to $\sqrt n$, then
\[
 \Pr\{\nu(A)\in\{h-1,h\}\}=1-O(2^{-c\sqrt n}).
\]
Along any sequence with $\min_{j\ge0}|n-j^2|\to\infty$, one has
$\Pr\{\nu(A)=\lfloor\sqrt n\rfloor\}\to1$.
\end{theorem}

The number $2^{n-r^2}$ has a counting interpretation. A fixed central principal block contains about a constant times $2^{-r^2}$ many $r$-dimensional kernel subspaces in expectation, and there are about $2^n$ choices of the block. Expectation alone leaves open whether these subspaces usually occur together on a small exceptional set of matrices. The second moment decides this issue.

For example, at $n=k^2+k$ the maximum equals $k$ with probability tending to one. At $n=k^2-1$ both $k-1$ and $k$ occur with probabilities bounded away from zero. The distinction depends on the distance of the \emph{dimension} $n$ from an integer square.

\section{Starting point and scope}
Vaas posed the typical-size problem for this exact binary group and also asked whether a linear lower bound could hold with overwhelming probability~\cite{vaas}. The ProveIt packet~\cite{prev} proved the finite upper bound $6\,2^{n-r(r-1)/2}$ for $1\le r\le\lfloor n/2\rfloor$. That establishes a square-root upper scale and rules out the proposed linear behavior. Its proof counts independent frames in a symplectic space and handles the all-ones vector separately. It explicitly leaves open the matching lower bound. Our estimate in Lemma~\ref{lem:fibre} retains all the Gram conditions when comparing witness probabilities, and Section~\ref{sec:lower} controls their simultaneous occurrence.

The larger archive~\cite{catalog} contains 61 research packets exploring consequences and method transfers from the shorter Lean proof of RH and Dirichlet GRH. The index records that they were deposited without independent checking there. The corresponding upstream theorem statements~\cite{lean} use Mathlib's zeta function and Dirichlet $L$-functions. They include nonvanishing to the right of $1/2$ for every Dirichlet character of positive modulus, with the principal pole excluded. One of the upstream reductions bounds a smoothed sum of von Mangoldt weights twisted by a quadratic character on a progression, then uses a one-sided criterion to obtain nonvanishing. Sutherland records independent verification of both editions at the pinned revision~\cite{suth}.

The requested RH article connects zero location to errors in prime counts~\cite{wiki}. The companion paper uses that link for primes sorted by multiplicative characters. The present paper follows the finite-field question identified by ProveIt. Its proof below is independent of the analytic theorem. Related papers treat nullity under principal pivot transforms~\cite{pivot}, distributions for one random matrix~\cite{rank}, and principal ranks in graph-state ensembles~\cite{graph}. The result here concerns the maximum over all principal blocks of one uniform dot-product isometry. Our novelty claim is relative to the explicit question and the inspected sources. The two papers in this archive present mathematical proofs with exact finite checks.

\section{A relative comparison for prescribed images}\label{sec:fibre}

Write $\mathbf j_m=(1,\ldots,1)^{\mathsf T}\in\mathbb F_2^m$. Call an
$m\times d$ matrix $Y$ regular if its columns are independent and
$\mathbf j_m$ lies outside their span. This extra condition matters in
characteristic two, since every orthogonal matrix fixes $\mathbf j_m$.

\begin{lemma}[Comparison on free rows]\label{lem:fibre}
Let $X$ be a regular $n\times d$ binary matrix, let $U$ be uniform in
$O(n,2)$ for the standard dot product, and let $Y$ be an $n\times d$
matrix of independent fair bits. If a positive-probability event $E$ for
the entries of an $n\times d$ matrix leaves $m\ge 2d+5$ specified rows
completely unrestricted, then
\[
 \left|\frac{\Pr(UX\in E)}{\Pr(Y\in E)}-1\right|
 \le 8\,2^{2d-m}.
\]
Here $E$ may impose arbitrary joint conditions on the remaining rows. The
bound is uniform in $X$, its Gram matrix, and the conditions.
\end{lemma}

\begin{proof}
First, the orbit of $X$ consists of all regular frames with its Gram
matrix $G=X^{\mathsf T}X$. To see this in odd dimension, the projection
$v\mapsto v+(v\cdot v)\mathbf j_n$ along $\mathbf j_n$ has symplectic
image $\mathbf j_n^\perp$ and is injective on the span of each regular
frame. Two frames with the same Gram matrix induce an isometry of the
projected subspaces. Their projected pairings equal
$v\cdot w+(v\cdot v)(w\cdot w)$. The symplectic form of Witt's extension
lemma extends this isometry to $\mathbf j_n^\perp$. Indeed, in each
projected subspace split off a nondegenerate symplectic summand and
choose a basis of its radical. In the orthogonal complement of the
summand, pair the first radical basis vector with a vector of pairing
one against it and zero against the remaining radical basis vectors.
Nondegeneracy gives such a vector. Split off this hyperbolic plane,
repeat, then complete the residual space by symplectic pairs. Matching
the resulting bases gives the extension, including when the projected
subspace is degenerate. Fixing $\mathbf j_n$ gives the required orthogonal
map. In even dimension adjoin one new unit vector $e$ and append $e$
to each frame, after embedding all its columns with last coordinate zero.
The augmented frames are regular, because
$\mathbf j_{n+1}=\mathbf j_n+e$. The odd-dimensional construction fixes
$e$, so it restricts to $e^\perp$. Consequently $UX$ has the same law as
$Y$ conditioned on regularity and $Y^{\mathsf T}Y=G$.

We next give local counting bounds. Let $Z$ be an $m\times d$ matrix of
independent fair bits, let $D=d(d+1)/2$, let $H$ be any symmetric
$d\times d$ binary matrix, and set $a=2^{2d-m}$. We claim that
\begin{align}
 2^D\Pr(Z^{\mathsf T}Z=H,\ Z\text{ regular}) &\ge 1-\tfrac43a,
 \label{eq:fibre-lower}\\
 2^D\Pr(Z^{\mathsf T}Z=H) &\le (1-4a)^{-1}.
 \label{eq:fibre-upper}
\end{align}
For the first inequality, expose the columns of $Z$. Given $i$ previous
regular columns, the $i$ scalar products and the prescribed norm of the
next column are $i+1$ independent linear constraints, because its norm
is its dot product with $\mathbf j_m$. Of the $2^{m-i-1}$ choices
satisfying these constraints, at most $2^{i+1}$ lie in the span of
$\mathbf j_m$ and the earlier columns. The remaining choices preserve
regularity. Multiply the resulting bounds for $0\le i<d$, then use
$\prod(1-t_i)\ge1-\sum t_i$ and
$\sum_{i=0}^{d-1}2^{2i+2-m}<4a/3$.

For the upper bound, a Fourier character on symmetric matrices,
evaluated at one row contribution $zz^{\mathsf T}$, has the form
$(-1)^{Q(z)}$, where $Q$ is a squarefree polynomial of degree at most
two and zero constant term. If its alternating polar form $B$ has
rank $2s$, the bias $b_Q=\mathbb E_z(-1)^{Q(z)}$ satisfies
\[
 b_Q^2=2^{-d}\sum_{h\in\operatorname{rad}B}(-1)^{Q(h)}.
\]
This follows by writing a pair of independent inputs as $z,z+h$ and
averaging first over $z$. Thus $b_Q$ vanishes unless $Q$ vanishes on
the radical, and otherwise $|b_Q|=2^{-s}$. For a fixed $B$ at most
$2^{2s}$ linear adjustments to $Q$ vanish on its radical. There are at
most $2^{2ds}$ choices for $B$ of rank $2s$, since such a form is a sum
of $s$ terms $uv^{\mathsf T}+vu^{\mathsf T}$ with $u,v\in\mathbb F_2^d$.
The only rank-zero character of nonzero bias is the trivial character.
By Fourier inversion and independence of the $m$ rows, the absolute
deviation of $2^D\Pr(Z^{\mathsf T}Z=H)$ from $1$ is at most
$\sum_{s\ge1}2^{(2d+2-m)s}=4a/(1-4a)$. This proves
\eqref{eq:fibre-upper}.

Condition on all the restricted rows of $Y$. They contribute a fixed
symmetric matrix to its Gram matrix. If the free-row matrix is regular,
the entire matrix is regular, as can be seen by restricting any proposed
linear dependence to the free rows. Thus
\eqref{eq:fibre-lower}--\eqref{eq:fibre-upper} sandwich the conditional
probability of a regular frame of Gram matrix $G$ between
$2^{-D}(1-4a/3)$ and $2^{-D}(1-4a)^{-1}$. The same bounds hold after
averaging with or without $E$. Bayes' rule and the orbit description
give the asserted ratio bound, since $a\le1/32$ and
$[(1-4a)(1-4a/3)]^{-1}-1\le8a$.
\end{proof}

\section{Complementation and the first moment}\label{sec:upper}
For a support $S$, write $V_S$ for its coordinate subspace. The kernel of a block is identified with $V_S\cap A^{-1}V_{S^c}$. Applying $A$ identifies this intersection with the corresponding kernel of $A^{-1}[S^c,S^c]$. Orthogonality gives $A^{-1}=A^{\mathsf T}$ and a square matrix and its transpose have the same nullity. Consequently,
\begin{equation}\label{eq:comp}
 \dim\ker A[S,S]=\dim\ker A[S^c,S^c].
\end{equation}
In particular the maximum is at most $\lfloor n/2\rfloor$. We can represent every complementary pair by a support of size $b\le n/2$, choosing one in case of a tie.

Let $U\subseteq V_S$ have dimension $r\ge1$. A basis matrix $X$ for $U$ has independent columns and its span omits the all-ones vector. It lies in the kernel of the block precisely when the rows of $AX$ in $S$ are zero. This event has probability $2^{-br}$ in an independent uniform $n\times r$ image matrix. For $n-b\ge2r+5$ Lemma~\ref{lem:fibre} gives the relative error $8\,2^{2r-(n-b)}$. There are
\[
 G(b,r)=\prod_{i=0}^{r-1}\frac{2^b-2^i}{2^r-2^i}
 \le c_\infty 2^{r(b-r)},\qquad
 c_\infty=\prod_{i\ge1}(1-2^{-i})^{-1}<4
\]
possible planes if $b\ge r$, by counting ordered bases. The bound on $c_\infty$ follows since the product of the first two factors $1-2^{-i}$ times $1-\sum_{i\ge3}2^{-i}$ equals $9/32$. For $b<r$ the event is empty. A union bound over the $2^{n-1}$ complementary pairs therefore gives, whenever $n/2\ge2r+5$,
\begin{equation}\label{eq:upper}
 \Pr(\nu(A)\ge r)\le 2^{n-1}c_\infty 2^{-r^2}
      (1+8\,2^{2r-n/2}).
\end{equation}
\section{The second moment over central supports}
\label{sec:lower}

Write $V_S=\{x\in\mathbf F_2^n:\operatorname{supp}x\subseteq S\}$.
An $r$-plane $U\le V_S$ witnesses nullity at least $r$ if
$(Ax)_i=0$ for every $x\in U$ and $i\in S$.
Let $W_S(A)$ be the number of such planes. It equals the Gaussian
coefficient $\genfrac{[}{]}{0pt}{}{\operatorname{null} A[S,S]}{r}_2$
and depends only on the nullity. We count
\[
 Z=\sum_{S\in\mathcal C} W_S(A),\qquad
 \mathcal C=\{S\subseteq[n]: n/3\le |S|\le 2n/3\}.
\]
For a uniform random subset the size is binomial with variance $n/4$.
Chebyshev's inequality gives $|\mathcal C|\ge(1-9/n)2^n$.
Constants in the following bounds are absolute. We work throughout with
$\tfrac12\sqrt n\le r\le 2\sqrt n$ and sufficiently large $n$.

\paragraph{One support.}
Use $A_0$ for the matrix with independent uniform binary entries.
Write $G(b,j)$ for the number of $j$-planes in $\mathbf F_2^b$,
with value $0$ for $j>b$. Set
$c_r=\prod_{i=1}^r(1-2^{-i})^{-1}$ and
$c_\infty=\prod_{i\ge1}(1-2^{-i})^{-1}<\infty$.
Each fixed plane for $S$ has witness probability $2^{-br}$, where
$b=|S|$, so
\[
 \mu_b:=\mathbb E W_S(A_0)
 =G(b,r)2^{-br}
 =c_r2^{-r^2}\prod_{i=0}^{r-1}(1-2^{i-b}).
\]
The final product differs from $1$ by at most $2^{r-b}$.
Lemma~\ref{lem:fibre} applies to each plane since its vectors vanish
off $S$ and the event ignores $n-b$ rows. Consequently, uniformly on
$\mathcal C$,
\begin{equation}\label{eq:onecentral}
 \mathbb E W_S(A)=(1+O(2^{-c n}))\,c_r2^{-r^2}
\end{equation}
for an absolute $c>0$.

\paragraph{Two supports in the independent model.}
Fix $S,R\in\mathcal C$, write $b=|S|$, $b'=|R|$,
$a=|S\cap R|$ and $D=b+b'-2a$. Draw $U\le V_S,V\le V_R$
independently and uniformly from their respective $r$-planes.
Write $i=\dim(U\cap V)$. Each row in $S\cap R$ carries $2r-i$
independent constraints, and each row in $S\triangle R$ carries
$r$. Their joint probability divided by $2^{-(b+b')r}$ is $2^{ai}$.
Thus
\[
 \frac{\mathbb E[W_S(A_0)W_R(A_0)]}{\mu_b\mu_{b'}}
 =\mathbb E_{U,V}2^{ai}.
\]
Count linear maps from $\mathbf F_2^a$ to $U\cap V$ by their image
dimension to obtain
\[
 2^{ai}=\sum_{j=0}^{\min(a,i)}G(i,j)\prod_{h=0}^{j-1}(2^a-2^h).
\]
Any common $j$-plane lies in $V_{S\cap R}$. The probability that a
given such plane is contained in $U$ is
$\prod_{h=0}^{j-1}(2^r-2^h)/(2^b-2^h)\le 2^{-j(b-r)}$
for $j\le r$, with the analogous bound for $V$.
These probabilities multiply. Using
$G(a,j)\le c_\infty 2^{j(a-j)}$ yields
\begin{equation}\label{eq:overlapmoment}
 0\le \mathbb E_{U,V}2^{ai}-1
 \le B_D,\qquad
 B_D:=c_\infty\sum_{j=1}^r 2^{j(2r-D-j)}.
\end{equation}

\paragraph{Complementation before the comparison.}
Fix $S,T\in\mathcal C$. Choose $R$ from $T,T^c$ to minimize
$D=|S\triangle R|$, resolving ties arbitrarily. Then $D\le n/2$,
$R\in\mathcal C$, and
\[
 n-|S\cup R|=n-\frac{|S|+|R|+D}{2}\ge n/12.
\]
Complementary principal blocks of an orthogonal matrix have equal
nullities, so $W_T(A)=W_R(A)$ pointwise. For any pair of planes on
$S,R$, choose a basis of their sum. Its span omits the all-ones vector
since $S\cup R$ omits a coordinate. The joint witness event constrains
the rows of the image of this basis and ignores at least $n/12$ rows.
Lemma~\ref{lem:fibre}, with at most $2r$ columns, therefore bounds its
probability by $1+O(2^{-c n})$ times the independent-entry value,
after decreasing the absolute constant $c$ if necessary. Indeed, a fair independent row applied to a matrix of full column rank is uniform on the whole space of output rows. Summing over
planes and using \eqref{eq:onecentral} and \eqref{eq:overlapmoment}, we get
\[
 \frac{\mathbb E Z^2}{(\mathbb E Z)^2}
 \le (1+O(2^{-c n}))\bigl(1+\mathbb E_{S,T} B_D\bigr),
\]
where $S,T$ are independent uniform elements of $\mathcal C$.
The complement identity was used only under the orthogonal law.

\paragraph{Averaging the overlap.}
For fixed $S$, the fraction of $T$ giving distance $D=d$ is at most
$2\binom nd/|\mathcal C|\ll 2^{-n}\binom nd$, since either $T$ or
its complement lies at Hamming distance $d$ from $S$.
Completing the square in \eqref{eq:overlapmoment} gives
$B_d\ll 2^{(r-d/2)^2}$, since the Gaussian sums
$\sum_{j\in\mathbf Z}2^{-(j-t)^2}$ are uniformly bounded in real $t$.
If $d<r$, then
\[
 \binom nd B_d\ll 2^{r^2+d\log_2n-rd+d^2/4}
       \le 2^{r^2-rd/2}
\]
up to an absolute multiplicative constant, for sufficiently large $n$.
Summing over $d<r$ contributes $O(2^{r^2-n})$ to $\mathbb E B_D$.
If $r\le d\le 3r$, the quadratic exponent in each summand of
\eqref{eq:overlapmoment} is at most $r^2/4$. The contribution of
these distances is $O(2^{-n+r^2/4+O(r\log n)})$.
Finally, for $d>3r$, summing the geometric majorant gives
$B_d\ll 2^{-r}$, irrespective of the probability of these distances.

Since $r\asymp\sqrt n$, the middle contribution is $O(2^{r^2-n})$
throughout the stated range: $3r^2/4$ dominates $O(r\log n)$.
Combining all three ranges gives
\begin{equation}\label{eq:lowerfailure}
 \frac{\operatorname{Var} Z}{(\mathbb E Z)^2}
 \ll 2^{r^2-n}+2^{-r},
 \qquad \tfrac12\sqrt n\le r\le2\sqrt n .
\end{equation}
The multiplicative error $O(2^{-c n})$ on the second moment can be
absorbed here even when $2^{r^2-n}$ is large. Chebyshev gives
$\Pr(\nu(A)<r)\le\Pr(Z=0)\le\operatorname{Var} Z/(\mathbb E Z)^2$.
Cauchy--Schwarz also gives
$\Pr(\nu(A)\ge r)\ge\Pr(Z>0)\ge(\mathbb E Z)^2/\mathbb E Z^2$.
These bounds give the asserted lower estimates on either side of the threshold.
\section{Consequences and the remaining critical window}
\begin{proof}[Completion of the proof of Theorem~\ref{thm:threshold}]
Equations \eqref{eq:upper} and \eqref{eq:lowerfailure} give \eqref{eq:fail}--\eqref{eq:hit}, the latter using the Cauchy--Schwarz lower bound for $\Pr(Z>0)$. Since $(h-1/2)^2\le n\le(h+1/2)^2$, both $n-(h-1)^2$ and $(h+1)^2-n$ are at least $h-1$. Apply \eqref{eq:fail} at $r=h-1$ and \eqref{eq:hit} at $r=h+1$. If $k=\lfloor\sqrt n\rfloor$ and $n$ recedes from the squares, the two quantities $n-k^2$ and $(k+1)^2-n$ both tend to infinity. Bounds at $k$ and $k+1$ now give the single-value conclusion.
\end{proof}

\begin{corollary}[Two values are sometimes necessary]
Along $n=k^2-1$ both $\Pr(\nu(A)=k-1)$ and $\Pr(\nu(A)=k)$ have strictly positive lower limits.
\end{corollary}
\begin{proof}
Theorem~\ref{thm:threshold} puts the maximum in $\{k-1,k\}$ with probability tending to one. Its lower hit bound at $r=k$ is at least $c/2$. By \eqref{eq:upper}, the upper hit probability is at most $c_\infty/4+o(1)<1$. Both assertions follow.
\end{proof}

\section{The leading constant above the threshold}
\label{sec:raretail}

Suppose $r/\sqrt n\to1$ and $r^2-n\to+\infty$. Then
\begin{equation}\label{eq:raretail}
 \Pr\bigl(\nu(A)\ge r\bigr)
 \ \sim\ \frac12\prod_{i=1}^\infty(1-2^{-i})^{-1}\,2^{n-r^2}.
\end{equation}
The factor $1/2$ is due to complementary supports giving the same event.

To prove this, set $\lambda=2^{n-r^2}$ and choose one member of
each complementary pair of central supports. Let $N$ count the
chosen supports with nullity at least $r$. For a block of nullity
$h$, counting flags of dimensions $r,r+1$ gives
\[
 G(h,r)(2^{h-r}-1)=G(h,r+1)(2^{r+1}-1).
\]
Thus, on the event $h\ge r+1$, its witness count $G(h,r)$ is at most
$(2^{r+1}-1)G(h,r+1)$. Formula~\eqref{eq:onecentral}, at both $r$
and $r+1$, shows that replacing the witness count by the indicator of
$h\ge r$ costs at most $O(2^{-r})$ of its mean. It follows that
$\mathbb E N\sim c_\infty\lambda/2$.

For the second factorial moment, use
$\mathbf 1_{h\ge r}\le G(h,r)$ at each of two distinct representatives
and repeat the pair comparison from Section~\ref{sec:lower}. Now $D=0$
is impossible. Each representative $S$ has at most $2\binom nd$ partners
at distance $D=d$. Summation of \eqref{eq:overlapmoment} for
$1\le d<r$ therefore contributes $O(\lambda^{-1}2^{-r/2})$ to the
normalized pair bound. The distances $r\le d\le3r$ contribute $o(1)$
because $r^2/n\to1$. Larger distances contribute $O(2^{-r})$.
Consequently
\[
 \mathbb E[N(N-1)]\ll \lambda^2+\lambda2^{-r/2}=o(\lambda).
\]
The first two Bonferroni inequalities give
$\Pr(N>0)=\mathbb E N+o(\lambda)$.
For all remaining supports, take the smaller member of each complement
pair, of size $b<n/3$. Supports of size $b<r$ cannot witness the
event. For each other support, the single-plane comparison on $n-b$
free rows gives $\mathbb E W_S(A)\ll 2^{-r^2}$.
There are $o(2^n)$ noncentral supports, so their contribution is
$o(\lambda)$ by a union bound. This proves \eqref{eq:raretail}.

\section{The full linear-deviation rate}
\begin{corollary}\label{cor:linear}
For every fixed $0<\alpha\le1/2$,
\[
 \lim_{n\to\infty}\frac{1}{n^2}\log_2\Pr\{\nu(A)\ge\lfloor\alpha n\rfloor\}=-\alpha^2 .
\]
\end{corollary}
\begin{proof}
Write $g(r)=r(r+1)/2$. First suppose $r\le(n-2)/2$. A basis $X$ of an $r$-plane supported on $S$, $|S|=b\le n/2$, is regular. In the full $n$-space the orbit of $X$ has at least $(2/3)2^{nr-g(r)}$ elements. Indeed the column-exposure proof of \eqref{eq:fibre-lower} gives the factor $1-(4/3)2^{2r-n}\ge2/3$ in this range, independently of the hypothesis used for the final ratio in Lemma~\ref{lem:fibre}. Within the $m=n-b$ coordinates outside $S$, there are at most $2^{mr-r(r-1)/2}$ independent frames of the prescribed Gram matrix. After $i$ columns, the prescribed pairings with them give $i$ independent linear restrictions on the next column in this nondegenerate $m$-space. Dropping norm and independence conditions only enlarges each successive count. Dividing by the full orbit size and summing over the planes and supports gives
\[
 \Pr\{\nu(A)\ge r\}\le \tfrac32 c_\infty 2^{n-r^2+r}.
\]
For the possible remaining threshold $r=\lfloor n/2\rfloor$, using $r-1$ on the right still gives $2^{-r^2+O(n)}$.

For a lower bound take one fixed $S$ of size $r$ and require $A[S,S]=0$. Let $X$ consist of its $r$ coordinate basis vectors, so $X^{\mathsf T}X=I_r$. Count choices of orthonormal images in the $m=n-r\ge r$ complementary coordinates. In a nondegenerate, nonalternating binary space of dimension $k$ the norm is a nonzero linear functional, so exactly $2^{k-1}$ vectors have norm one. Its representing vector $c$ satisfies $v\cdot v=v\cdot c$ for all $v$. Upon splitting off a norm-one vector $y$, the nondegenerate complement is alternating only if $c\in\langle y\rangle$, hence $c=y$. Thus at most one choice is forbidden if a positive-dimensional nonalternating complement is required. In dimension one the unique choice is allowed. The product of the fractions of allowed choices, over the successive dimensions $k=m,m-1,\ldots,m-r+1$, is at least $\prod_{k\ge2}(1-2^{1-k})=1/c_\infty$. This gives at least $2^{mr-g(r)}/c_\infty$ frames. They all omit the global all-ones vector from their span. In the full space there are at most $2^{nr-g(r)}$ orthonormal $r$-frames: at each step the residual norm-one fiber is empty or has exactly half the residual space. The orbit characterization therefore gives probability at least $2^{-r^2}/c_\infty$. Divide logarithms of the upper and lower bounds by $n^2$ and let $r=\lfloor\alpha n\rfloor$.
\end{proof}

\paragraph{Exact finite checks.}
The accompanying Python program enumerates the orthogonal groups in dimensions 2 through 5, compares frame orbits with their complete Gram fibers for frames of up to three columns, and tests the complementary-block identity. In a separate enumeration of all $65{,}536$ independent-entry binary $4\times4$ matrices, it counts one-dimensional kernel witnesses on all six two-coordinate supports. It finds $\mathbb E Z=9/2$ and $\mathbb E Z^2=201/8$, agreeing with the row-wise and overlap counts. Further checks treat the Gram bounds themselves for every symmetric target through three columns. The estimates for growing $n$ follow from the proofs.

\paragraph{Further research questions.}
For fixed $t\in\mathbb Z$, does $\Pr\{\nu(A)\ge k\}$ converge along $n=k^2+t$, and what is its limit? A candidate counting argument must account for complementary supports giving exactly the same nullity. For $0<\alpha<1/2$, can one determine the term of order $n$ in the logarithm of the linear-deviation probability after subtracting $-\lfloor\alpha n\rfloor^2$? The proof of Corollary~\ref{cor:linear} permits a gap of that order.

\clearpage
\begin{thebibliography}{9}
\bibitem{vaas} Vaas, \href{https://math.stackexchange.com/questions/5119604/principal-submatrices-of-random-binary-orthogonal-matrices}{Principal submatrices of random binary orthogonal matrices}, Mathematics Stack Exchange, question 5119604 (14 January 2026).
\bibitem{prev} ProveIt research archive, \href{https://github.com/VladimirReshetnikov/ProveIt/blob/02dbe7d5754f1e0b6879185e60302813738de756/Analysis/GeneralizedRiemannHypothesis/Research/grh_method_transfers_orthogonal_jets_2026-11-26/sections/binary.tex}{A tail bound for every principal binary orthogonal kernel}, in \emph{Beyond Dirichlet GRH: orthogonal kernels and sharp $p$-adic recovery}, commit \texttt{02dbe7d5} (2026).
\bibitem{catalog} \href{https://github.com/VladimirReshetnikov/ProveIt/tree/02dbe7d5754f1e0b6879185e60302813738de756/Analysis/GeneralizedRiemannHypothesis}{ProveIt, Analysis/GeneralizedRiemannHypothesis}. Directory index and provenance notes, same revision as \cite{prev}.
\bibitem{lean} L. Sernau et al., GRH for Dirichlet $L$-functions, in Lean 4 (2026). Pinned \href{https://github.com/lsernau/grh-lean/blob/72fead102f663d5c7a407964a13dc0cf9d977207/grh/src/Statements.lean}{statements}, \href{https://github.com/lsernau/grh-lean/blob/72fead102f663d5c7a407964a13dc0cf9d977207/grh/src/Main.lean}{exports}, and \href{https://github.com/lsernau/grh-lean/blob/72fead102f663d5c7a407964a13dc0cf9d977207/grh/src/New/Reduction/OneSidedBound.lean}{one-sided reduction}.
\bibitem{suth} A. V. Sutherland, \href{https://math.mit.edu/~drew/shortgrh/}{Verification files for the shorter Lean proof of RH and of GRH for Dirichlet $L$-functions} (26 November 2026).
\bibitem{wiki} \href{https://en.wikipedia.org/wiki/Riemann_hypothesis}{Riemann hypothesis}, Wikipedia, consulted 26 November 2026, especially the formulation and distribution-of-primes sections.
\bibitem{pivot} R. Brijder and H. J. Hoogeboom, \href{https://arxiv.org/abs/0912.0878}{Nullity invariance for pivot and the interlace polynomial}, \emph{Linear Algebra Appl.} 435 (2011), 277--288.
\bibitem{rank} J. Fulman and L. Goldstein, \href{https://arxiv.org/abs/1211.0504}{Stein's method and the rank distribution of random matrices over finite fields}, \emph{Ann. Probab.} 43 (2015), 1274--1314.
\bibitem{graph} S. Ghosh, D. Hangleiter and J. Helsen, \href{https://doi.org/10.1103/52xz-3hpc}{Random regular graph states are complex at almost any depth}, \emph{PRX Quantum} 6 (2025), 040344.
\end{thebibliography}
\end{document}
